\section{Thuật ngữ và ký hiệu}
\label{sec:notation}

Trước hết, chúng ta có thể đọc qua một số thuật ngữ và ký hiệu sẽ được dùng trong bài viết.

\subsection{Bảng thuật ngữ}

\begin{center}
  \small
  \begin{tabular}{@{}>{\raggedright\arraybackslash}p{0.28\textwidth}
    >{\raggedright\arraybackslash}p{0.62\textwidth}@{}}
    \toprule
    \textbf{Thuật ngữ} & \textbf{Ý nghĩa}                                                                                        \\
    \midrule
    LTSF
                       & Long-term Time Series Forecasting: dự báo chuỗi thời gian với horizon dài.                              \\
    \addlinespace
    Linear / NLinear / DLinear
                       & Ba baseline tuyến tính trong LTSF-Linear: hồi quy thuần, chuẩn hóa bằng giá trị cuối,
    và phân rã bằng trung bình trượt rồi dự báo hai nhánh.                                                                       \\
    \addlinespace
    Decomposition (phân rã)
                       & Tách chuỗi thành xu hướng, mùa (và phần dư), rồi mô hình hóa từng phần.                                 \\
    \addlinespace
    Normalization (chuẩn hóa)
                       & Ở đây chủ yếu là recentering/inversing bằng giá trị lookback cuối (NLinear), không phải StandardScaler. \\
    \addlinespace
    Data characteristics
                       & Đặc trưng định lượng của chuỗi (độ mạnh xu hướng/mùa, mức không dừng) dùng để giải thích
    vì sao một inductive bias thắng.                                                                                             \\
    \addlinespace
    Lookback \(L\)
                       & Độ dài cửa sổ quá khứ đưa vào mô hình.                                                                  \\
    \addlinespace
    Horizon / pred\_len \(T\)
                       & Độ dài đoạn tương lai cần dự báo.                                                                       \\
    \addlinespace
    VIC
                       & Chuỗi giá đóng cửa hàng ngày của cổ phiếu Vingroup (HOSE), dùng \(\log(\mathrm{close})\).
                         Đây là dữ liệu tài chính thêm vào để kiểm tra xem kết luận từ các bộ LTSF chuẩn
                         có còn đúng trên chuỗi giá cổ phiếu (nhiễu, dễ đổi trạng thái thị trường) hay không. \\
    \addlinespace
    STL
                       & Seasonal-Trend decomposition bằng Loess: tách chuỗi thành xu hướng, mùa, phần dư.                       \\
    \addlinespace
    Channel-independent
                       & Mỗi kênh biến được dự báo riêng; trọng số \(W\) dùng chung giữa các kênh.                               \\
    \bottomrule
  \end{tabular}
\end{center}

\subsection{Bảng ký hiệu toán học}

\begin{center}
  \small
  \begin{tabular}{@{}>{\raggedright\arraybackslash}p{0.22\textwidth}
    >{\raggedright\arraybackslash}p{0.68\textwidth}@{}}
    \toprule
    \textbf{Ký hiệu} & \textbf{Ý nghĩa}                                                                    \\
    \midrule
    \(L\)
                     & Độ dài lookback (số bước quá khứ).                                                  \\
    \addlinespace
    \(T\)
                     & Độ dài horizon (số bước dự báo). Trong STL, \(T\) cũng ký hiệu thành phần xu hướng;
    ngữ cảnh công thức sẽ làm rõ.                                                                          \\
    \addlinespace
    \(C\)
                     & Số kênh (biến) của chuỗi đa biến.                                                   \\
    \addlinespace
    \(x\), \(\hat{y}\)
                     & Cửa sổ đầu vào và dự báo đầu ra.                                                    \\
    \addlinespace
    \(l\)
                     & Giá trị lookback cuối cùng; NLinear trừ rồi cộng lại \(l\).                         \\
    \addlinespace
    \(m\)
                     & Kích thước kernel trung bình trượt của DLinear.                                     \\
    \addlinespace
    \(S\), \(R\)
                     & Thành phần mùa và phần dư từ STL.                                                   \\
    \addlinespace
    \(F_T\), \(F_S\)
                     & Độ mạnh xu hướng và độ mạnh mùa (Hyndman / FPP3), nằm trong \([0,1]\).              \\
    \addlinespace
    \(\mathrm{MSE}\), \(\mathrm{MAE}\)
                     & Mean Squared Error và Mean Absolute Error.                                          \\
    \addlinespace
    \(\Delta\)
                     & \(\mathrm{MSE}(\mathrm{DLinear}) - \mathrm{MSE}(\mathrm{NLinear})\);
    âm nghĩa là DLinear tốt hơn.                                                                           \\
    \bottomrule
  \end{tabular}
\end{center}
